\scanpage{71}
\textcircled{2} $d'=$ tous les ``petits'' diagrammes, \inserted{$d''=\varnothing$.} Pas de conditions d'exactitude.

\uline{Pb 1} $I$ une petite catégorie, on a pour $C\in\mathrm{Ob}\,\Delta$
\[
\underline{\mathrm{Hom}}(I,C)\simeq\underline{\mathrm{Hom}}(I^\circ,
\underset{\substack{\text{\unsure{avec|sans}}\\\varinjlim\ \text{quelc.}}}{C^\circ})^{\circ}
\simeq\underline{\mathrm{Hom}}_{\Delta^\circ}(\widehat{I^\circ},C^\circ)^\circ
\simeq\underline{\mathrm{Hom}}_{\Delta}(\widehat{I^\circ}{}^\circ,C)
\]
donc $I$ a la $\Delta$-enveloppe $\widetilde I=\underline{\mathrm{Hom}}(I,\unsure{C|e})^\circ$ (qui est l'opposé d'un topos). Notons que $I\longrightarrow\widetilde I$ définit les objets de $I$ comme famille \uline{cogénératrice} (donc $I^\circ\longrightarrow\widetilde I^\circ$ comme famille génératrice).

\uline{Pb 2} $C\in\mathrm{Ob}\,\Delta$, $\Sigma\subset\mathrm{Fl}\,C$ [$\Sigma$ petit]. Rendre inversibles les éléments de $\Sigma$ ? \deleted{Supposons que} Or $C\Sigma^{-1}$ (au sens de Cat) \inserted{est défini, montrer que c'est une solution,}. On \uline{peut supposer} que $\Sigma$ est stable par composition, par chgt de base, contient les isom, et \inserted{tel} que $vu\in\Sigma$, $v\in\Sigma\Longrightarrow u\in\Sigma$. Mais alors $\Sigma$ donne lieu à un ``calcul de fractions \unsure{à g|: g}'', et il est connu que \inserted{dans $C\Sigma^{-1}$} les \unsure{dp|types} pour les petites $\varprojlim$ qui existent dans $\unsure{\mathscr C|\mathbf C|C}$ existent dans $C\Sigma^{-1}$, et $C\longrightarrow C\Sigma^{-1}$ y commute.
% slanted
\margin{$\Sigma$ ne sera plus petit après saturation !}
Il reste à examiner les produits \inserted{(infinis)}. Pour ceci, on peut supposer aussi que $\Sigma$ est stable par produits \uline{infinis} de morphismes. Mais alors on voit tout de suite que les produits de $C$ sont des produits de $C\Sigma^{-1}$. OK !
% slanted
\margin{Mais il faut \deleted{Si $C$ est} vérifier que la catégorie obtenue est une $\mathcal U$-catégorie.}

Notons que si $I$ est une famille ``cogénératrice \unsure{pour mon strict|pour mon stricts|par mono strict}'' dans \deleted{Pour tout} $C$ [pour tt $X\in\mathrm{Ob}\,C$, $X\xrightarrow{\sim}\varprojlim_{X^{\prime}/C}X'$], alors il en est encore de \unsure{même|$=$} dans $C\Sigma^{-1}$. \deleted{Donc on voit}

Supposons qu'on ait un \inserted{\unsure{premier|pseudo-}} syst. inductif filtrant $(R_\alpha)$ dans $\Delta$, avec $I\longrightarrow R_\alpha$ cogénérateur. Soit $R=\underset{\alpha}{\varinjlim}_{\Delta}R_\alpha$. Alors $I\longrightarrow R$ est-il cogénérateur ? Soit $R'=\underset{\alpha}{\varinjlim}_{\mathrm{Cat}}R_\alpha$. \deleted{Dans $R'$, les $\varprojlim$ \unsure{prennent|proviennent} de syst.} Donc $R$ s'obtient à partir de $R'$ en rendant inversibles certaines flèches, donc $\mathrm{Ob}\,R^{\unsure{\prime\prime|\prime}}=\bigcup_{\alpha}\mathrm{Im}\,R_\alpha$. De ceci on déduit aussitôt ce qu'on a dit.
% slanted
\margin{\uline{Mais attention} le fait d'être strictement cogénérateur \uline{n'est} \uline{pas} transitif}

Il faut aussi étudier le cas où on ``rajoute des flèches'' à $R$ muni de $I$, d'où $R\longrightarrow R'$. Ce \deleted{\unsure{dernier}} \inserted{$\Delta$-foncteur} est ess. surjectif, et on gagne encore !
% archive sheet number: 70
